% ============================================================
%  AUTOCORRELACIÓN — ECONOMETRÍA
%  Capítulo 12 · Wooldridge & Gujarati
%  Estilo académico B&W (NYT / The Economist)
%  Clase: 17 de marzo de 2026
% ============================================================

\documentclass[11pt, letterpaper]{article}

% ── ENCODING & IDIOMA ──────────────────────────────────────
\usepackage[utf8]{inputenc}
\usepackage[T1]{fontenc}
\usepackage[spanish,es-nodecimaldot]{babel}

% ── TIPOGRAFÍA ─────────────────────────────────────────────
\usepackage{lmodern}
\usepackage{microtype}
\usepackage{setspace}
\onehalfspacing

% ── GEOMETRÍA ──────────────────────────────────────────────
\usepackage[
  top=2.5cm, bottom=2.8cm,
  left=3.2cm, right=2.8cm,
  headheight=14pt
]{geometry}

% ── MATEMÁTICAS ────────────────────────────────────────────
\usepackage{amsmath,amssymb,amsthm}
\usepackage{mathtools}
\usepackage{bm}

% ── GRÁFICOS ───────────────────────────────────────────────
\usepackage{tikz}
\usetikzlibrary{arrows.meta,positioning,calc,shapes.geometric,
                decorations.pathmorphing,patterns}
\usepackage{pgfplots}
\pgfplotsset{compat=1.18}
\usepackage{graphicx,float}

% ── TABLAS ─────────────────────────────────────────────────
\usepackage{booktabs,array,tabularx,multirow,longtable}

% ── COLORES — ESCALA DE GRISES PURA ────────────────────────
\usepackage{xcolor}
\definecolor{Black}{RGB}{0,0,0}
\definecolor{DarkGray}{RGB}{40,40,40}
\definecolor{MedGray}{RGB}{120,120,120}
\definecolor{LightGray}{RGB}{220,220,220}
\definecolor{VeryLight}{RGB}{246,246,246}

% ── CAJAS (B&W) ────────────────────────────────────────────
\usepackage{tcolorbox}
\tcbuselibrary{skins,breakable,theorems}

\tcbset{
  defstyle/.style={
    enhanced, breakable,
    colback=VeryLight, colframe=Black,
    fonttitle=\bfseries\small\scshape, title=#1,
    left=8pt, right=8pt, top=5pt, bottom=5pt,
    borderline west={2.5pt}{0pt}{Black},
    sharp corners
  },
  thmstyle/.style={
    enhanced, breakable,
    colback=white, colframe=DarkGray,
    fonttitle=\bfseries\small\scshape, title=#1,
    left=8pt, right=8pt, top=5pt, bottom=5pt,
    borderline west={4pt}{0pt}{DarkGray},
    sharp corners
  },
  notestyle/.style={
    enhanced, breakable,
    colback=LightGray!40, colframe=MedGray,
    fonttitle=\bfseries\small\scshape, title=#1,
    left=8pt, right=8pt, top=5pt, bottom=5pt,
    borderline west={2pt}{0pt}{MedGray},
    sharp corners
  },
  formulastyle/.style={
    enhanced, breakable,
    colback=white, colframe=Black,
    fonttitle=\bfseries\normalsize\scshape, title=#1,
    left=10pt, right=10pt, top=8pt, bottom=8pt,
    sharp corners,
    frame style={line width=1.2pt}
  }
}

\newtcolorbox{definicion}[1]{defstyle={#1}}
\newtcolorbox{teorema}[1]{thmstyle={#1}}
\newtcolorbox{nota}[1]{notestyle={#1}}
\newtcolorbox{formulabox}[1]{formulastyle={#1}}

% ── CABECERAS ──────────────────────────────────────────────
\usepackage{fancyhdr}
\pagestyle{fancy}
\fancyhf{}
\fancyhead[L]{\small\textit{Econometría \textperiodcentered\ Autocorrelación}}
\fancyhead[R]{\small\textit{Cap.\,12 --- Wooldridge \& Gujarati}}
\fancyfoot[C]{\small\thepage}
\renewcommand{\headrulewidth}{0.6pt}

% ── HIPERVÍNCULOS (B&W) ────────────────────────────────────
\usepackage[colorlinks=true,
            linkcolor=black,
            citecolor=black,
            urlcolor=black]{hyperref}
\usepackage{cleveref}

% ── LISTAS ─────────────────────────────────────────────────
\usepackage{enumitem}
\setlist[enumerate]{leftmargin=*,itemsep=3pt,topsep=4pt}
\setlist[itemize]{leftmargin=*,itemsep=3pt,topsep=4pt}

% ── FORMATO DE SECCIONES estilo NYT ────────────────────────
\usepackage{titlesec}
\titleformat{\section}
  {\large\bfseries\scshape}{\thesection.}{0.6em}{}
  [\vspace{-4pt}\rule{\linewidth}{0.6pt}\vspace{2pt}]
\titleformat{\subsection}
  {\normalsize\bfseries}{\thesubsection.}{0.5em}{}
\titleformat{\subsubsection}
  {\normalsize\itshape}{\thesubsubsection.}{0.5em}{}

% ── TÍTULO ─────────────────────────────────────────────────
\title{%
  \rule{\linewidth}{1.4pt}\\[6pt]
  {\Large\bfseries\scshape Econometr\'ia --- Notas de Clase}\\[4pt]
  {\large Cap\'itulo 12: Autocorrelaci\'on en los Errores}\\[3pt]
  {\normalsize\textit{Wooldridge, Cap.\,12 \quad\&\quad Gujarati \& Porter, Cap.\,12}}\\[6pt]
  \rule{\linewidth}{0.6pt}
}
\author{\normalsize Apuntes de clase en la UPTC por Emanuel Quintana Silva \quad\textperiodcentered\quad 17 de marzo de 2026}
\date{}

% ════════════════════════════════════════════════════════════
\begin{document}
\maketitle
\thispagestyle{fancy}

\begin{abstract}
\noindent
Estos apuntes desarrollan de forma exhaustiva el problema de la
\textbf{autocorrelaci\'on} (o correlaci\'on serial) en los errores del modelo de
regresi\'on lineal cl\'asico. Se cubren: definici\'on y naturaleza del problema,
las ocho causas principales, el proceso AR(1) y sus propiedades, las
consecuencias sobre los estimadores MCO, los m\'etodos gr\'aficos y param\'etricos
de detecci\'on (Rachas, Durbin--Watson, Breusch--Godfrey) con su derivaci\'on
alg\'ebrica completa, los procedimientos remediales (GLS, FGLS, Newey--West/HAC)
y un \'arbol de decisi\'on integrador.
\end{abstract}

\newpage

\tableofcontents

\newpage

% ============================================================
\section{Contexto: Supuestos del MCRL sobre los Errores}
% ============================================================

El Modelo Cl\'asico de Regresi\'on Lineal (MCRL) impone sobre el vector de
perturbaciones $\mathbf{u}$ la estructura de covarianza escalar:
\[
  E[\mathbf{u}\mathbf{u}^\top] = \sigma^2 \mathbf{I}_n
\]
lo que implica dos supuestos simult\'aneos:
\begin{enumerate}
  \item \textbf{Homocedasticidad:} $\operatorname{Var}(u_i) = \sigma^2$ para todo $i$.
  \item \textbf{No autocorrelaci\'on:}
    $\operatorname{Cov}(u_i, u_j) = E[u_i u_j] = 0$ para $i \neq j$.
\end{enumerate}

\begin{definicion}{Supuesto de No Autocorrelaci\'on (Gauss--Markov)}
\[
  \operatorname{Cov}(u_i,\, u_j) \;=\; E\!\left[u_i\, u_j\right] \;=\; 0,
  \qquad i \neq j
\]
Para el modelo $Y_i = \beta_1 + \beta_2 X_i + u_i$: el error de cualquier
observaci\'on \emph{no} est\'a relacionado con el error de ninguna otra
observaci\'on, ni en el \textbf{tiempo} ni en el \textbf{espacio}.
Si este supuesto se viola: $E[\mathbf{u}\mathbf{u}^\top] = \sigma^2\Omega$
donde $\Omega \neq \mathbf{I}$.
\end{definicion}

\begin{figure}[H]
\centering
\begin{tikzpicture}[>=Stealth, thick]
  \draw[->] (0,0) -- (5.8,0) node[right] {$X_i$};
  \draw[->] (0,0) -- (0,4.8) node[above] {$Y_i$};
  \draw[very thick] (0.5,0.8) -- (5.2,4.1)
        node[right,font=\small] {$\hat{Y}_i = \hat\beta_1+\hat\beta_2 X_i$};
  \foreach \x/\y/\lab in {
    1.0/2.3/{u_m},
    1.8/1.1/{u_s},
    2.6/2.9/{u_i},
    3.4/3.5/{u_l},
    4.2/2.7/{u_r}}{
    \filldraw[Black] (\x,\y) circle (2.5pt);
    \node[font=\tiny,above right] at (\x,\y) {$\lab$};
    \pgfmathsetmacro{\yfit}{0.74 + 0.638*\x}
    \draw[dashed,MedGray,thin] (\x,\y) -- (\x,\yfit);
  }
\end{tikzpicture}
\caption{Supuesto cumplido: residuos $u_i$ dispersos aleatoriamente,
sin patr\'on sistem\'atico en el tiempo ni el espacio.}
\end{figure}

% ============================================================
\section{Naturaleza del Problema}
% ============================================================

\begin{definicion}{Definici\'on --- Gujarati \& Porter (2010, p.\,412)}
La \textbf{autocorrelaci\'on} es la correlaci\'on entre miembros de series de
observaciones ordenadas en el tiempo (series de tiempo) o en el espacio
(datos de corte transversal). Ocurre cuando $E(u_i u_j)\neq 0$ para $i\neq j$.
\end{definicion}

Cuando existe autocorrelaci\'on, la matriz de covarianza de los errores adopta:
\[
  E[\mathbf{u}\mathbf{u}^\top] = \sigma^2
  \begin{pmatrix}
    1          & \rho_1    & \rho_2    & \cdots & \rho_{n-1} \\
    \rho_1     & 1         & \rho_1    & \cdots & \rho_{n-2} \\
    \vdots     & \ddots    & \ddots    & \ddots & \vdots     \\
    \rho_{n-1} & \cdots    & \cdots    & \cdots & 1
  \end{pmatrix}
  = \sigma^2 \Omega
\]
donde $\rho_k = \operatorname{Corr}(u_t,\, u_{t-k})$ es la autocorrelaci\'on en el rezago $k$.

\subsection{Distinci\'on Terminol\'ogica}

\begin{center}
\begin{tabular}{lp{9.5cm}}
\toprule
\textbf{T\'ermino} & \textbf{Definici\'on (Tintner)} \\
\midrule
\textbf{Autocorrelaci\'on} &
  Correlaci\'on de una serie \emph{consigo misma} con rezago:
  $\operatorname{Corr}(u_t,\, u_{t-k})$, $k \geq 1$. \\[5pt]
\textbf{Correlaci\'on serial} &
  Correlaci\'on entre dos series \emph{distintas} desplazadas en el tiempo.
  En pr\'actica econom\'etrica ambos t\'erminos se usan indistintamente. \\
\bottomrule
\end{tabular}
\end{center}

\subsection{Tipos de Datos y Manifestaci\'on}

\begin{enumerate}
  \item \textbf{Series de tiempo} --- forma m\'as frecuente; errores del
    periodo $t$ correlacionados con los de $t-1,\,t-2,\ldots$
  \item \textbf{Datos de corte transversal} --- autocorrelaci\'on
    \emph{espacial}: unidades geogr\'aficas vecinas presentan errores similares.
  \item \textbf{Datos de panel} --- combinaci\'on de dependencia temporal
    y espacial (datos agrupados).
\end{enumerate}

% ============================================================
\section{Causas de la Autocorrelaci\'on}
% ============================================================

\begin{enumerate}[label=\textbf{\arabic*.}]

  \item \textbf{Inercia (Momentum macroecon\'omico).}
  Las series econ\'omicas (PIB, empleo, inflaci\'on, consumo) exhiben ciclos de
  expansi\'on y contracci\'on. En una fase de recuperaci\'on, cada valor sucesivo
  suele ser mayor que el anterior, creando dependencia sistem\'atica entre
  observaciones contiguas.

  \item \textbf{Sesgo de especificaci\'on: variables excluidas.}
  Si se omite una variable relevante \emph{que s\'i est\'a autocorrelacionada},
  su influencia se absorbe en $u_t$.

  \medskip
  \noindent\textit{Ejemplo (apuntes):}
  \begin{align}
    \text{Especificado:} &\quad
      I_t = \beta_1 + \beta_2\,\text{Ingreso}_t + \beta_3\,\text{Impuestos}_t
            + \underbrace{u_t}_{\ni\,\alpha_4\,r_t} \label{eq:omit}\\
    \text{Correcto:} &\quad
      I_t = \alpha_1 + \alpha_2\,\text{Ingreso}_t + \alpha_3\,\text{Impuestos}_t
            + \alpha_4\,r_t + v_t \notag
  \end{align}
  Al omitir la tasa de inter\'es $r_t$, el error de \eqref{eq:omit}
  hereda su autocorrelaci\'on.

  \item \textbf{Forma funcional incorrecta.}
  Ajustar un modelo lineal a una relaci\'on curvilínea produce residuos con
  patr\'on sistem\'atico (sobreestimaci\'on en los extremos, subestimaci\'on
  en el centro).

  \item \textbf{Omisi\'on de variables rezagadas.}
  Ignorar la din\'amica del modelo (efectos retardados). Ejemplo agr\'icola:
  \begin{align}
    \text{Incorrecto:} &\quad S_t = \beta_1 + \beta_2 P_t - \beta_3 C_t + u_t \\
    \text{Correcto:}   &\quad S_t = \beta_1 + \beta_2 P_{t-1} - \beta_3 C_{t-1} + u_t
  \end{align}
  donde $S_t=$ oferta agr\'icola, $P_t=$ precio, $C_t=$ costos de producci\'on.

  \item \textbf{Fen\'omeno de la telara\~na (\textit{cobweb}).}
  En mercados con rezago entre producci\'on y venta, la oferta actual responde
  al precio del periodo anterior, generando un patr\'on c\'iclico de
  sobreproducci\'on/subproducci\'on que se refleja en los errores del modelo.

  \item \textbf{Manipulaci\'on de datos.}
  Promediar, interpolar o extrapolar datos para cambiar la frecuencia
  (mensual $\to$ trimestral) introduce suavizado artificial que correlaciona
  observaciones adyacentes.

  \item \textbf{Transformaciones de datos.}
  Aplicar primeras diferencias ($\Delta y_t = y_t - y_{t-1}$) para corregir
  colinealidad puede \emph{inducir} autocorrelaci\'on en $\Delta u_t$ incluso
  si $u_t$ era independiente.

  \item \textbf{No estacionariedad.}
  Si las variables del modelo no son estacionarias, los errores reflejan
  esa inestabilidad mediante autocorrelaci\'on persistente ($\rho \to 1$).

\end{enumerate}

\begin{nota}{Autocorrelaci\'on Pura vs.\ Inducida}
La autocorrelaci\'on \textbf{pura} surge de la naturaleza del fen\'omeno
(inercia, cobweb). La \textbf{inducida} es un artefacto de errores del
investigador: especificaci\'on incorrecta, omisi\'on de variables o
transformaciones inadecuadas. Antes de aplicar remedios es fundamental
diagnosticar cu\'al tipo est\'a presente.
\end{nota}

% ============================================================
\section{El Esquema Autorregresivo AR(1)}
% ============================================================

\begin{definicion}{Proceso AR(1) --- Esquema de Primer Orden}
\[
  \boxed{u_t \;=\; \rho\, u_{t-1} \;+\; \varepsilon_t},
  \qquad -1 < \rho < 1
\]
\begin{itemize}
  \item $\rho$ --- coeficiente de autocorrelaci\'on de primer orden
    (\emph{poblacional}); $\mathrm{AR}(1) \equiv \rho$.
  \item $\varepsilon_t$ --- ruido blanco puro:
    (1)~$E[\varepsilon_t]=0$; \;
    (2)~$\operatorname{Var}(\varepsilon_t)=\sigma_\varepsilon^2$; \;
    (3)~$\operatorname{Cov}(\varepsilon_t,\varepsilon_s)=0$, $t\neq s$.
\end{itemize}
\end{definicion}

\subsection{Propiedades del Proceso AR(1) Estacionario ($|\rho|<1$)}

\begin{align}
  E[u_t] &= 0 \label{eq:ar1mean}\\[4pt]
  \operatorname{Var}(u_t) &= \frac{\sigma_\varepsilon^2}{1-\rho^2}
    \label{eq:ar1var}\\[4pt]
  \operatorname{Cov}(u_t,\, u_{t-k}) &= \rho^k \cdot
    \frac{\sigma_\varepsilon^2}{1-\rho^2}
    \label{eq:ar1cov}\\[4pt]
  \operatorname{Corr}(u_t,\, u_{t-k}) &= \rho^k
    \label{eq:ar1corr}
\end{align}

\noindent\textit{Derivaci\'on de \eqref{eq:ar1var}:}
Tomando varianza de $u_t = \rho u_{t-1}+\varepsilon_t$:
\[
  \sigma_u^2 = \rho^2\sigma_u^2 + \sigma_\varepsilon^2
  \;\Rightarrow\;
  \sigma_u^2(1-\rho^2) = \sigma_\varepsilon^2
  \;\Rightarrow\;
  \sigma_u^2 = \frac{\sigma_\varepsilon^2}{1-\rho^2}
\]

\begin{nota}{Interpretaci\'on de $\rho$ y su relaci\'on con $d$}
\begin{center}
\begin{tabular}{ccl}
\toprule
\textbf{Valor de $\rho$} & $d \approx 2(1-\rho)$ & \textbf{Interpretaci\'on}\\
\midrule
$\rho = 0$    & $d = 2$ & Sin autocorrelaci\'on\\
$\rho \to 1$  & $d \to 0$ & Autocorr.\ positiva perfecta\\
$\rho \to -1$ & $d \to 4$ & Autocorr.\ negativa perfecta\\
\bottomrule
\end{tabular}
\end{center}
Rango: $-1 \leq \rho \leq 1$ y $0 \leq d \leq 4$.
\end{nota}

\subsection{Esquemas de Orden Superior}

\begin{itemize}
  \item \textbf{AR($p$):}\quad
    $u_t = \rho_1 u_{t-1} + \rho_2 u_{t-2} + \cdots + \rho_p u_{t-p} + \varepsilon_t$
  \item \textbf{MA(1):}\quad $u_t = \varepsilon_t + \lambda\varepsilon_{t-1}$
  \item \textbf{ARMA($p,q$):} combinaci\'on autorregresiva y de media m\'ovil.
\end{itemize}
La prueba Durbin--Watson s\'olo detecta AR(1). Para \'ordenes superiores se
requiere Breusch--Godfrey.

% ============================================================
\section{Consecuencias de Estimar con MCO bajo Autocorrelaci\'on}
% ============================================================

\begin{teorema}{Teorema --- Propiedades de MCO bajo AR(1)}
En presencia de autocorrelaci\'on AR(1), el estimador MCO $\hat{\bm\beta}$:
\begin{enumerate}
  \item \textbf{Linealidad e insesgadez:}
    $E[\hat{\bm\beta}] = \bm\beta$ \quad (se preserva: depende de
    $E[\mathbf{u}]=\mathbf{0}$, no de $E[\mathbf{u}\mathbf{u}^\top]$).
  \item \textbf{Consistencia:}
    $\hat{\bm\beta} \xrightarrow{p} \bm\beta$ \quad (se preserva).
  \item \textbf{Eficiencia (BLUE) perdida:}
    MCO \textbf{ya no es el mejor estimador lineal insesgado}.
    El estimador MCG/GLS tiene menor varianza.
\end{enumerate}
\end{teorema}

\subsection{Distorsi\'on de la Varianza del Estimador}

F\'ormula habitual de MCO (incorrecta bajo AR(1)):
\[
  \widehat{\operatorname{Var}}(\hat\beta_2)_{\text{MCO}} = \frac{\hat\sigma^2}{\sum x_t^2}
\]

Varianza \emph{verdadera} bajo AR(1):
\[
  \operatorname{Var}(\hat\beta_2)_{\text{AR1}}
  = \frac{\sigma^2}{\sum x_t^2}
    \left[
      1 + 2\rho\,\frac{\sum x_t x_{t-1}}{\sum x_t^2}
        + 2\rho^2\,\frac{\sum x_t x_{t-2}}{\sum x_t^2} + \cdots
    \right]
\]

Si $X_t$ tambi\'en sigue AR(1) con coeficiente $r$:
\[
  \boxed{
    \operatorname{Var}(\hat\beta_2)_{\text{AR1}}
    \approx
    \operatorname{Var}(\hat\beta_2)_{\text{MCO}}
    \cdot \frac{1 + r\rho}{1 - r\rho}
  }
\]
Cuando $r > 0$ y $\rho > 0$ (frecuente en macroeconom\'ia), este factor
$>1$: MCO \textbf{subestima} la varianza real.

\subsection{Estimador $\hat\sigma^2$ Sesgado}

\[
  \hat\sigma^2 = \frac{\sum_{t=1}^n \hat{u}_t^2}{n-k}
\]
es un estimador \textbf{sesgado hacia abajo} de $\sigma^2$ bajo
autocorrelaci\'on positiva, agravando la subestimaci\'on de los errores
est\'andar.

\subsection{Cuadro Completo de Consecuencias}

\begin{center}
{\small
\begin{tabular}{llp{6.5cm}}
\toprule
\textbf{Propiedad} & \textbf{Estado} & \textbf{Consecuencia pr\'actica}\\
\midrule
Insesgadez           & Se mantiene &
  $E[\hat\beta]=\beta$; no hay sesgo sistem\'atico.\\
Consistencia         & Se mantiene &
  Con $n\to\infty$, $\hat\beta\to\beta$.\\
Eficiencia (BLUE)    & Se pierde   &
  Existe GLS con menor varianza.\\
$\hat\sigma^2$       & Sesgado $\downarrow$ &
  Falsa sensaci\'on de precisi\'on.\\
Errores est\'andar    & Subestimados &
  Estad\'isticos $t$ y $F$ artificialmente altos.\\
Pruebas $t$ y $F$    & Inv\'alidas  &
  Error Tipo~I excesivo; variables declaradas\\
                     &             &
  significativas cuando no lo son.\\
$R^2$                & Sobreestimado &
  Falsa sensaci\'on de buen ajuste.\\
Intervalos de conf.  & Demasiado estrechos &
  Inferencia sobre $\beta$ no es confiable.\\
\bottomrule
\end{tabular}
}
\end{center}

% ============================================================
\section{M\'etodos de Detecci\'on}
% ============================================================

\subsection{M\'etodo Gr\'afico}

Estimar el modelo por MCO, obtener $\hat{u}_t$ y graficarlos contra
(a)~el tiempo $t$, o (b)~el residuo rezagado $\hat{u}_{t-1}$.

\begin{figure}[H]
\centering
%── Fila (a): autocorrelaci\'on positiva ───────────────────────
\begin{minipage}[t]{0.47\textwidth}
\centering
\begin{tikzpicture}[scale=0.92]
  \draw[->](-0.2,0)--(5.7,0) node[right,font=\small]{Tiempo};
  \draw[->](0,-1.7)--(0,1.9) node[above,font=\small]{$u_t$};
  \node[font=\small,left] at (0,0){$0$};
  \foreach \i in {0,...,33}{
    \pgfmathsetmacro{\xx}{0.165*\i+0.15}
    \pgfmathsetmacro{\yy}{1.30*sin(deg(1.12*\xx))}
    \filldraw[Black](\xx,\yy)circle(1.8pt);
  }
\end{tikzpicture}
\end{minipage}
\hfill
\begin{minipage}[t]{0.47\textwidth}
\centering
\begin{tikzpicture}[scale=0.92]
  \draw[->](-1.9,0)--(2.1,0) node[right,font=\small]{$u_{t-1}$};
  \draw[->](0,-1.9)--(0,2.1) node[above,font=\small]{$u_t$};
  \node[font=\small,below left] at (0,0){$0$};
  \foreach \i in {0,...,39}{
    \pgfmathsetmacro{\t}{2*3.14159*\i/40}
    \pgfmathsetmacro{\u}{1.55*cos(deg(\t))}
    \pgfmathsetmacro{\v}{0.32*sin(deg(\t))+0.08*sin(deg(3*\t))}
    \pgfmathsetmacro{\px}{0.7071*\u - 0.7071*\v}
    \pgfmathsetmacro{\py}{0.7071*\u + 0.7071*\v}
    \filldraw[Black](\px,\py)circle(1.8pt);
  }
\end{tikzpicture}
\end{minipage}
\begin{center}\small\textit{(a) Autocorrelaci\'on positiva:
onda suave / nube el\'iptica con pendiente $+45^{\circ}$}\end{center}

\bigskip

%── Fila (b): autocorrelaci\'on negativa ──────────────────────
\begin{minipage}[t]{0.47\textwidth}
\centering
\begin{tikzpicture}[scale=0.92]
  \draw[->](-0.2,0)--(5.7,0) node[right,font=\small]{Tiempo};
  \draw[->](0,-1.7)--(0,1.9) node[above,font=\small]{$u_t$};
  \node[font=\small,left] at (0,0){$0$};
  \foreach \i/\y in {
    0/1.1,1/-1.2,2/0.9,3/-1.3,4/1.0,5/-0.8,6/1.2,
    7/-1.1,8/0.9,9/-1.2,10/1.0,11/-0.9,12/1.3,
    13/-1.0,14/0.8,15/-1.2,16/1.1}{
    \pgfmathsetmacro{\xx}{0.33*\i+0.18}
    \filldraw[Black](\xx,\y)circle(1.8pt);
  }
  \draw[Black!50,thin]
    (0.18,1.1)--(0.51,-1.2)--(0.84,0.9)--(1.17,-1.3)
    --(1.50,1.0)--(1.83,-0.8)--(2.16,1.2)--(2.49,-1.1)
    --(2.82,0.9)--(3.15,-1.2)--(3.48,1.0)--(3.81,-0.9)
    --(4.14,1.3)--(4.47,-1.0)--(4.80,0.8)--(5.13,-1.2)
    --(5.46,1.1);
\end{tikzpicture}
\end{minipage}
\hfill
\begin{minipage}[t]{0.47\textwidth}
\centering
\begin{tikzpicture}[scale=0.92]
  \draw[->](-1.9,0)--(2.1,0) node[right,font=\small]{$u_{t-1}$};
  \draw[->](0,-1.9)--(0,2.1) node[above,font=\small]{$u_t$};
  \node[font=\small,below left] at (0,0){$0$};
  \foreach \i in {0,...,39}{
    \pgfmathsetmacro{\t}{2*3.14159*\i/40}
    \pgfmathsetmacro{\u}{1.55*cos(deg(\t))}
    \pgfmathsetmacro{\v}{0.32*sin(deg(\t))+0.08*sin(deg(3*\t))}
    \pgfmathsetmacro{\px}{ 0.7071*\u + 0.7071*\v}
    \pgfmathsetmacro{\py}{-0.7071*\u + 0.7071*\v}
    \filldraw[Black](\px,\py)circle(1.8pt);
  }
\end{tikzpicture}
\end{minipage}
\begin{center}\small\textit{(b) Autocorrelaci\'on negativa:
zigzag r\'apido / nube el\'iptica con pendiente $-45^{\circ}$}\end{center}

\caption{Figura 12.3 (Gujarati \& Porter, 2010, p.\,416).
Columna izquierda: $u_t$ contra el tiempo.
Columna derecha: $u_t$ contra $u_{t-1}$.}
\end{figure}

\begin{figure}[H]
\centering
\begin{tikzpicture}[scale=0.92]
  \draw[->](-0.2,0)--(5.7,0) node[right,font=\small]{Tiempo};
  \draw[->](0,-1.7)--(0,1.9) node[above,font=\small]{$u_t$};
  \node[font=\small,left] at (0,0){$0$};
  \draw[dashed,MedGray,thin](0,0)--(5.5,0);
  \foreach \x/\y in {
    0.3/0.9, 0.8/-0.7, 1.3/1.2, 1.9/-1.0, 2.4/0.4,
    2.9/1.1, 3.5/-0.5, 4.0/0.8, 4.5/-1.1, 5.0/0.3}{
    \filldraw[Black](\x,\y)circle(1.8pt);
  }
\end{tikzpicture}
\caption{Sin autocorrelaci\'on: residuos $\hat{u}_t$ dispersos aleatoriamente
alrededor de cero, sin ning\'un patr\'on sistem\'atico.}
\end{figure}

% ============================================================
\subsection{Prueba de las Rachas (Geary / Wald--Wolfowitz)}
% ============================================================

\subsubsection{Hip\'otesis}
\begin{align*}
  H_0 &: \text{Los } \hat{u}_t \text{ son aleatorios (no est\'an
          correlacionados en el tiempo).}\\
  H_a &: \text{Los } \hat{u}_t \text{ no son aleatorios ---
          existe correlaci\'on serial.}
\end{align*}

\subsubsection{Procedimiento}
\begin{enumerate}
  \item Estimar el modelo por MCO; obtener $\hat{u}_1,\ldots,\hat{u}_n$.
  \item Asignar signos ($+$ o $-$) a cada residuo.
  \item Identificar las \textbf{rachas}: sucesiones ininterrumpidas del
    mismo signo.
\end{enumerate}

\begin{definicion}{Definici\'on --- Racha}
Una \textbf{racha} es una subsecuencia contigua de observaciones con el mismo
s\'imbolo (signo). Ejemplo con $R=5$ rachas:
\[
  \underbrace{+,+,+}_{R_1},\;
  \underbrace{-,-}_{R_2},\;
  \underbrace{+}_{R_3},\;
  \underbrace{-,-,-,-}_{R_4},\;
  \underbrace{+,+}_{R_5}
\]
\end{definicion}

\noindent Sea $N_1=$ n\'umero de $+$, $N_2=$ n\'umero de $-$, $N=N_1+N_2$.

\subsubsection{Distribuci\'on Bajo $H_0$}

\begin{align}
  E(R) &= \frac{2 N_1 N_2}{N} + 1 \label{eq:Emean}\\[6pt]
  \sigma_R^2 &= \frac{2 N_1 N_2 (2 N_1 N_2 - N)}{N^2 (N-1)}
  \label{eq:Rvar}
\end{align}

Para $N$ grande ($N_1>10$, $N_2>10$):
$Z = (R - E(R))/\sigma_R \sim \mathcal{N}(0,1)$.

\subsubsection{Regla de Decisi\'on}

\[
  \Pr\!\left[\,E(R) - z_{\alpha/2}\,\sigma_R
             \leq R \leq
             E(R) + z_{\alpha/2}\,\sigma_R\,\right] = 1 - \alpha
\]

\begin{itemize}
  \item $R$ \textbf{dentro} del intervalo $\Rightarrow$
    \textbf{no rechazar} $H_0$.
  \item $R$ \textbf{fuera} del intervalo $\Rightarrow$
    \textbf{rechazar} $H_0$: los errores son interdependientes.
\end{itemize}

\subsubsection{Ejemplo Num\'erico}

\begin{center}
\begin{tabular}{ccc|ccc}
\toprule
$t$ & $\hat{u}_t$ & Signo &
$t$ & $\hat{u}_t$ & Signo \\
\midrule
1 & $+1.25$ & $+$ & 6 & $-0.63$ & $-$ \\
2 & $+2.53$ & $+$ & 7 & $-0.81$ & $-$ \\
3 & $-0.97$ & $-$ & 8 & $+4.10$ & $+$ \\
4 & $+1.58$ & $+$ & 9 & $+4.30$ & $+$ \\
5 & $-0.64$ & $-$ & & & \\
\bottomrule
\end{tabular}
\end{center}

\noindent Secuencia: $\underbrace{+,+}_{R_1},\underbrace{-}_{R_2},
\underbrace{+}_{R_3},\underbrace{-,-,-}_{R_4},\underbrace{+,+}_{R_5}$
$\;\Rightarrow\;$ $R=5$, $N_1=5$, $N_2=4$, $N=9$.
\[
  E(R) = \frac{2(5)(4)}{9}+1 \approx 5.44; \qquad
  \sigma_R \approx 1.38
\]
Al 95\%: intervalo $[2.74,\;8.14]$. Como $R=5 \in [2.74,\;8.14]$,
\textbf{no se rechaza $H_0$}.

% ============================================================
\subsection{Prueba $d$ de Durbin--Watson}
% ============================================================

\subsubsection{Hip\'otesis}
\begin{align*}
  H_0 &: \rho = 0 \quad\text{(no hay autocorrelaci\'on AR(1))}\\
  H_a &: \rho \neq 0 \quad\text{(hay autocorrelaci\'on de primer orden)}
\end{align*}

\subsubsection{Seis Supuestos Requeridos}
\begin{enumerate}
  \item El modelo incluye \textbf{intercepto}.
  \item Las variables explicativas son \textbf{no estoc\'asticas}
    (fijas en muestras repetidas).
  \item Las perturbaciones siguen AR(1):
    $u_t = \rho u_{t-1} + \varepsilon_t$.
  \item Los errores $\varepsilon_t$ se distribuyen \textbf{normalmente}.
  \item El modelo \textbf{no incluye} $Y_{t-1}$ como regresor.
  \item \textbf{No hay datos faltantes}.
\end{enumerate}

\subsubsection{Estad\'istico}

\[
  \boxed{
    d \;=\; \frac{\displaystyle\sum_{t=2}^{n}(\hat{u}_t - \hat{u}_{t-1})^2}
                 {\displaystyle\sum_{t=1}^{n}\hat{u}_t^2}
  }
\]

\subsubsection{Derivaci\'on Alg\'ebrica Completa}

Expandiendo el numerador:
\begin{align}
  d &= \frac{\sum \hat{u}_t^2
             - 2\sum \hat{u}_t\hat{u}_{t-1}
             + \sum \hat{u}_{t-1}^2}
            {\sum \hat{u}_t^2}
  \notag\\[6pt]
    &\approx
    \frac{2\sum \hat{u}_t^2 - 2\sum \hat{u}_t\hat{u}_{t-1}}
         {\sum \hat{u}_t^2}
  \tag{$\sum_{t=2}^n\hat u_t^2 \approx \sum_{t=2}^n\hat u_{t-1}^2$}
  \\[6pt]
    &= 2\!\left(1 - \underbrace{\frac{\sum\hat{u}_t\hat{u}_{t-1}}{\sum\hat{u}_t^2}}_{=\,\hat\rho}\right)
  \notag
\end{align}

\noindent Relaci\'on fundamental:
\[
  \boxed{d \;\approx\; 2(1 - \hat\rho)}, \qquad
  \hat\rho = \frac{\sum_{t=2}^n \hat{u}_t\hat{u}_{t-1}}{\sum_{t=1}^n\hat{u}_t^2},
  \qquad 0 \leq d \leq 4
\]

\subsubsection{Regiones de Decisi\'on}

\begin{figure}[H]
\centering
\begin{tikzpicture}[>=Stealth, scale=1.05]
  \draw[very thick] (0,0) -- (10,0);
  \foreach \x/\lab in {
    0/$0$, 2.0/$d_L$, 3.0/$d_U$,
    5.0/$2$, 7.0/$4{-}d_U$, 8.0/$4{-}d_L$, 10/$4$}{
    \draw (\x,0.12) -- (\x,-0.12);
    \node[below,font=\small] at (\x,-0.18) {$\lab$};
  }
  \fill[pattern=north east lines,pattern color=Black!55]
        (0,0.25) rectangle (2.0,0.55);
  \fill[LightGray](2.0,0.25) rectangle (3.0,0.55);
  \fill[white,draw=Black,thin](3.0,0.25) rectangle (7.0,0.55);
  \fill[LightGray](7.0,0.25) rectangle (8.0,0.55);
  \fill[pattern=north west lines,pattern color=Black!55]
        (8.0,0.25) rectangle (10,0.55);
  \node[font=\tiny,align=center] at (1.0,0.75)
    {Rechazar $H_0$\\autocorr.\ positiva};
  \node[font=\tiny,align=center] at (2.5,0.75){Inconcl.};
  \node[font=\tiny,align=center] at (5.0,0.75)
    {No rechazar $H_0$\\$(\rho = 0)$};
  \node[font=\tiny,align=center] at (7.5,0.75){Inconcl.};
  \node[font=\tiny,align=center] at (9.0,0.75)
    {Rechazar $H_0$\\autocorr.\ negativa};
\end{tikzpicture}
\caption{Regiones de decisi\'on del estad\'istico Durbin--Watson.}
\end{figure}

\begin{center}
{\small
\begin{tabular}{p{3.0cm}p{2.0cm}p{2.2cm}p{4.5cm}}
\toprule
\textbf{Rango de $d$} & $\hat\rho$ aprox. & \textbf{Decisi\'on} & \textbf{Interpretaci\'on}\\
\midrule
$0 < d < d_L$               & $\approx 1$  & Rechazar $H_0$ & Autocorr.\ positiva \\
$d_L \leq d \leq d_U$       & ---          & Inconcluso     & Zona de indeterminaci\'on \\
$d_U < d < 4{-}d_U$         & $\approx 0$  & No rechazar $H_0$ & Sin autocorrelaci\'on \\
$4{-}d_U \leq d \leq 4{-}d_L$ & ---        & Inconcluso     & Zona de indeterminaci\'on \\
$4{-}d_L < d < 4$            & $\approx -1$ & Rechazar $H_0$ & Autocorr.\ negativa \\
\bottomrule
\end{tabular}
}
\end{center}

\begin{nota}{Limitaciones de Durbin--Watson}
\begin{enumerate}
  \item S\'olo detecta autocorrelaci\'on de \textbf{primer orden} AR(1).
  \item \textbf{Inv\'alida} si el modelo incluye $Y_{t-1}$ como regresor.
  \item \textbf{Zona de indeterminaci\'on} entre $d_L$ y $d_U$.
  \item Los valores cr\'iticos dependen de $n$, $k$ y $\alpha$.
  \item Requiere errores normalmente distribuidos.
\end{enumerate}
\end{nota}

% ============================================================
\subsection{Estad\'istico $h$ de Durbin (Modelos Autorregresivos)}
% ============================================================

Cuando el modelo contiene $Y_{t-1}$ como regresor, $d$ est\'a sesgado.
Se usa el \textbf{estad\'istico $h$}:

\[
  h = \hat\rho \sqrt{\frac{n}{1 - n\,\widehat{\operatorname{Var}}(\hat\alpha)}}
  \;\xrightarrow{d}\; \mathcal{N}(0,1) \quad\text{bajo }H_0
\]

donde $\hat\alpha$ es el coeficiente MCO de $Y_{t-1}$.

\begin{nota}{Condici\'on de Validez}
$h$ no est\'a definido si $n\cdot\widehat{\operatorname{Var}}(\hat\alpha)\geq 1$.
En tal caso usar directamente Breusch--Godfrey.
\end{nota}

% ============================================================
\subsection{Prueba de Breusch--Godfrey (LM)}
% ============================================================

La prueba BG supera todas las limitaciones de Durbin--Watson:
\begin{enumerate}
  \item Detecta autocorrelaci\'on de \textbf{orden $p$ arbitrario}.
  \item V\'alida con regresores \textbf{estoc\'asticos}, incluyendo $Y_{t-1}$.
  \item No requiere normalidad de $\varepsilon_t$.
\end{enumerate}

\subsubsection{Procedimiento}
\begin{enumerate}
  \item Estimar el modelo original por MCO; obtener $\hat{u}_t$.
  \item Estimar la regresi\'on auxiliar:
  \[
    \hat{u}_t = \alpha_0 + \alpha_1 X_{1t} + \cdots + \alpha_k X_{kt}
                + \gamma_1 \hat{u}_{t-1} + \cdots + \gamma_p \hat{u}_{t-p}
                + \eta_t
  \]
  \item El estad\'istico LM:
  \[
    \text{LM} = (n-p)\,R^2 \;\xrightarrow{d}\; \chi^2(p)
    \quad\text{bajo }H_0:\,\gamma_1=\cdots=\gamma_p=0
  \]
  \item Versi\'on $F$ (muestras peque\~nas):
  \[
    F = \frac{R^2/p}{(1-R^2)/(n-k-p)} \sim F(p,\, n-k-p)
  \]
\end{enumerate}

% ============================================================
\section{Medidas Remediales}
% ============================================================

\subsection{M\'inimos Cuadrados Generalizados Exactos (MCG/GLS)}

\textbf{Caso: $\rho$ conocido.}
Transformaci\'on de cuasi-diferencias:
\begin{align*}
  Y_t^* &= Y_t - \rho\,Y_{t-1} \\
  X_t^* &= X_t - \rho\,X_{t-1}
\end{align*}

Modelo transformado ($t=2,\ldots,n$):
\[
  Y_t^* = \beta_1(1-\rho) + \beta_2 X_t^* + \varepsilon_t
\]
$\varepsilon_t$ es ruido blanco: se cumplen todos los supuestos del MCRL.

\subsubsection{Transformaci\'on de Prais--Winsten (Primera Observaci\'on)}

La primera observaci\'on no puede cuasi-diferenciarse. Se recupera con:
\[
  Y_1^* = \sqrt{1-\rho^2}\cdot Y_1, \qquad
  X_1^* = \sqrt{1-\rho^2}\cdot X_1
\]
En muestras peque\~nas, descartar la primera observaci\'on (Cochrane--Orcutt)
genera p\'erdida significativa de eficiencia.

\subsection{MCG Factibles (FGLS): $\rho$ Desconocido}

\subsubsection{M\'etodo de Cochrane--Orcutt (CO)}
\begin{enumerate}
  \item Estimar por MCO; obtener $\hat{u}_t$.
  \item Regresar $\hat{u}_t$ sobre $\hat{u}_{t-1}$ para estimar $\hat\rho$:
  \[
    \hat\rho = \frac{\sum_{t=2}^n \hat{u}_t\hat{u}_{t-1}}{\sum_{t=1}^n \hat{u}_t^2}
  \]
  \item Transformar: $\tilde{Y}_t = Y_t - \hat\rho\,Y_{t-1}$, $\tilde{X}_t = X_t - \hat\rho\,X_{t-1}$.
  \item Re-estimar por MCO con las variables transformadas.
  \item Actualizar $\hat\rho$ con los nuevos residuos y repetir hasta convergencia.
\end{enumerate}

\begin{nota}{CO vs.\ Prais--Winsten}
Cochrane--Orcutt \textbf{descarta} la primera observaci\'on en cada iteraci\'on
(costoso cuando $n<50$). Prais--Winsten la retiene aplicando $\sqrt{1-\hat\rho^2}$
y es preferible en la pr\'actica.
\end{nota}

\subsection{Errores Est\'andar Robustos --- Newey--West (HAC)}

En muestras \textbf{grandes}, se mantienen los coeficientes MCO pero
se corrige la inferencia con la matriz de varianza--covarianza HAC:
\[
  \widehat{\operatorname{Var}}(\hat{\bm\beta})_{\text{NW}}
  = (\mathbf{X}'\mathbf{X})^{-1}
    \hat{\mathbf{S}}
    (\mathbf{X}'\mathbf{X})^{-1}
\]
donde $\hat{\mathbf{S}} = \hat\Gamma_0 +
\sum_{j=1}^{m} w_j(\hat\Gamma_j+\hat\Gamma_j')$,
con pesos de Bartlett $w_j = 1 - \frac{j}{m+1}$
y truncamiento $m \approx \lfloor 4(n/100)^{2/9}\rfloor$.

\medskip\noindent\textbf{Propiedades:}
\begin{itemize}
  \item Coeficientes id\'enticos a MCO; s\'olo se corrigen los errores est\'andar.
  \item Consistente bajo heterocedasticidad y autocorrelaci\'on de cualquier
    orden hasta rezago $m$.
  \item No requiere especificar la forma de $\Omega$.
\end{itemize}

\newpage

% ============================================================
\section*{Hoja de F\'ormulas --- Referencia R\'apida}
\addcontentsline{toc}{section}{Hoja de F\'ormulas}
% ============================================================

\begin{formulabox}{F\'ormulas Esenciales --- Autocorrelaci\'on}
\begin{alignat}{3}
  &\text{Supuesto cl\'asico:}    &\quad&
    E(u_i u_j) = 0,\;\; i\neq j \tag{A1}\\[3pt]
  &\text{Proceso AR(1):}       &\quad&
    u_t = \rho\, u_{t-1} + \varepsilon_t,\;\; |\rho|<1 \tag{A2}\\[3pt]
  &\text{Varianza AR(1):}      &\quad&
    \operatorname{Var}(u_t) = \dfrac{\sigma_\varepsilon^2}{1-\rho^2} \tag{A3}\\[5pt]
  &\text{Covarianza AR(1):}    &\quad&
    \operatorname{Cov}(u_t,u_{t-k}) = \rho^k\dfrac{\sigma_\varepsilon^2}{1-\rho^2} \tag{A4}\\[5pt]
  &\text{Estad\'istico DW:}      &\quad&
    d = \dfrac{\sum_{t=2}^n(\hat u_t-\hat u_{t-1})^2}{\sum_{t=1}^n\hat u_t^2}
      \approx 2(1-\hat\rho) \tag{A5}\\[5pt]
  &\text{Estimador }\hat\rho: &\quad&
    \hat\rho = \dfrac{\sum \hat u_t\hat u_{t-1}}{\sum\hat u_t^2} \tag{A6}\\[5pt]
  &\text{Rango de }d: &\quad&
    0 \leq d \leq 4 \tag{A7}\\[3pt]
  &\text{Media de rachas:}     &\quad&
    E(R) = \dfrac{2N_1 N_2}{N}+1 \tag{A8}\\[5pt]
  &\text{Varianza de rachas:}  &\quad&
    \sigma_R^2 = \dfrac{2N_1 N_2(2N_1 N_2 - N)}{N^2(N-1)} \tag{A9}\\[5pt]
  &\text{Estad\'istico }h: &\quad&
    h = \hat\rho\sqrt{\dfrac{n}{1-n\,\widehat{\operatorname{Var}}(\hat\alpha)}} \tag{A10}\\[5pt]
  &\text{Prueba BG (LM):}      &\quad&
    \mathrm{LM} = (n-p)R^2 \sim \chi^2(p) \tag{A11}\\[3pt]
  &\text{Transf.\ GLS:}  &\quad&
    Y_t^* = Y_t - \rho Y_{t-1};\;\; X_t^* = X_t - \rho X_{t-1} \tag{A12}\\[3pt]
  &\text{Prais--Winsten }(t=1): &\quad&
    Y_1^* = \sqrt{1-\rho^2}\,Y_1;\;\; X_1^* = \sqrt{1-\rho^2}\,X_1 \tag{A13}\\[3pt]
  &\text{Var.\ verdadera MCO:} &\quad&
    \operatorname{Var}(\hat\beta_2)_{\text{AR1}}
    \approx \operatorname{Var}(\hat\beta_2)_{\text{MCO}}
    \cdot\dfrac{1+r\rho}{1-r\rho} \tag{A14}
\end{alignat}
\end{formulabox}

\bigskip

\begin{center}
{\small
\begin{tabular}{p{2.8cm}p{2.5cm}p{3.0cm}p{4.5cm}}
\toprule
\textbf{Prueba} & \textbf{Orden} & \textbf{Condici\'on} & \textbf{Limitaci\'on principal}\\
\midrule
Gr\'afico       & Cualquiera & Residuos MCO & Subjetivo\\
Rachas (Geary) & Cualquiera & S\'olo signos & No cuantifica magnitud\\
Durbin--Watson & AR(1)      & No $Y_{t-1}$ & Zona inconclusa; s\'olo orden\,1\\
Durbin $h$     & AR(1)      & Modelo AR    & Requiere $n\cdot\widehat{\text{Var}}<1$\\
Breusch--Godfrey & AR($p$)  & Residuos MCO & Elecci\'on de $p$\\
\bottomrule
\end{tabular}
}
\end{center}

\vspace{0.8cm}
\noindent\rule{\linewidth}{0.6pt}
\begin{center}
\small
Gujarati, D.\,N.\ \& Porter, D.\,C.\ (2010).
\textit{Econometr\'ia}, 5.\textsuperscript{a} ed., Cap.\,12. McGraw-Hill.\\
Wooldridge, J.\,M.\ (2016).
\textit{Introductory Econometrics}, 6th ed., Cap.\,12. Cengage.
\end{center}

\end{document}
